% !TeX program = xelatex
% =============================================================================
%  DỰ ÁN TỐT NGHIỆP — BÁO CÁO CHUYÊN ĐỀ
%  Đề tài: Phân tích Hiệu suất và Dự báo Sản lượng Hệ thống Điện Mặt Trời
%  Nhóm: The Outliers
%  Trường: Cao đẳng FPT Polytechnic — Chuyên ngành Xử lý Dữ liệu
% =============================================================================

\documentclass[12pt, a4paper]{report}

% ======================== PACKAGES CƠ BẢN & FONT ============================
\usepackage{geometry}
\geometry{a4paper, top=2.5cm, bottom=2.5cm, left=3cm, right=2cm, headheight=33pt, footskip=1.1cm}

% --- FONT CHO XELATEX (HỖ TRỢ TIẾNG VIỆT 100% - FORMAL, TỐI GIẢN, HIỆN ĐẠI) ---
\usepackage{fontspec}
\setmainfont{Times New Roman}
\setmonofont{Consolas}

% ======================== MÀU SẮC & CODE LISTINGS ===========================
\usepackage{xcolor}
\usepackage{listings}
% Định nghĩa bảng màu Tableau trích xuất từ main_fixed_01.twb & tableau_theme.json
\definecolor{tblNavy}{HTML}{4E79A7}
\definecolor{tblOrange}{HTML}{F28E2B}
\definecolor{tblGreen}{HTML}{59A14F}
\definecolor{tblRed}{HTML}{E15759}
\definecolor{tblTeal}{HTML}{76B7B2}
\definecolor{tblYellow}{HTML}{EDC948}
\definecolor{tblPurple}{HTML}{B07AA1}
\definecolor{tblGray}{HTML}{79706E}
\definecolor{tblSeqWarm}{HTML}{F28E2B}

% Định nghĩa màu sắc cho Code Block
\definecolor{mutedgreen}{rgb}{0.2,0.5,0.3} 
\definecolor{mutedblue}{rgb}{0.15,0.35,0.6} 
\definecolor{mutedgray}{rgb}{0.45,0.45,0.45} 
\definecolor{mutedpurple}{rgb}{0.45,0.2,0.55} 
\definecolor{softbackground}{rgb}{0.98,0.98,0.98} 
\definecolor{bordergray}{rgb}{0.85,0.85,0.85} 
\definecolor{fptorange}{RGB}{245,130,32}

\lstdefinestyle{pythonstyle}{
    backgroundcolor   = \color{softbackground},
    commentstyle      = \color{mutedgreen}\itshape,
    keywordstyle      = \color{mutedblue}\bfseries,
    numberstyle       = \tiny\color{mutedgray},
    stringstyle       = \color{mutedpurple},
    basicstyle        = \ttfamily\footnotesize\color{black},
    breakatwhitespace = false,
    breaklines        = true,
    columns           = fullflexible,
    captionpos        = t, 
    keepspaces        = true,
    numbers           = left,
    numbersep         = 5pt,
    showspaces        = false,
    showstringspaces  = false,
    showtabs          = false,
    tabsize           = 4,
    frame             = single,
    rulecolor         = \color{bordergray}
}
\lstset{style=pythonstyle}

% Ngăn chặn tự động gạch nối từ (hyphenation) cuối dòng
\hyphenpenalty=10000
\exhyphenpenalty=10000
\tolerance=10000
\emergencystretch=3em

% ======================== TRANG TRÍ & ĐỊNH DẠNG =============================
\usepackage{fancyhdr}
\usepackage{lastpage}
\usepackage{amsmath, amssymb}
\usepackage{enumitem}
\usepackage{array}
\newcolumntype{L}[1]{>{\raggedright\arraybackslash}p{#1}}
\newcolumntype{C}[1]{>{\centering\arraybackslash}p{#1}}
\newcolumntype{R}[1]{>{\raggedleft\arraybackslash}p{#1}}
\usepackage{booktabs}
\usepackage{tabularx}
\usepackage{longtable}
\usepackage{multirow}
\usepackage{graphicx}
\graphicspath{{images/}{figures/}{./}}
\usepackage{float}
\usepackage{caption}
\usepackage{subcaption}
\usepackage{tikz}
\usepackage[titles]{tocloft}

% VIỆT HÓA TOÀN BỘ CÁC TIÊU ĐỀ HỆ THỐNG
\renewcommand{\figurename}{Hình}
\renewcommand{\tablename}{Bảng}
\renewcommand{\contentsname}{Mục Lục}
\renewcommand{\listfigurename}{Danh Mục Hình Vẽ}
\renewcommand{\listtablename}{Danh Mục Bảng Biểu}
\providecommand{\tep}[1]{\texttt{#1}}

\usepackage{indentfirst}
\setlength{\parindent}{1cm}
\setlength{\parskip}{6pt}

\usepackage{setspace}
\onehalfspacing
\setlength{\headheight}{33pt}

% Hyperlinks & Bookmarks
\usepackage{hyperref}
\hypersetup{
    colorlinks  = true,
    linkcolor   = black,
    citecolor   = blue,
    filecolor   = magenta,
    urlcolor    = blue,
    pdftitle    = {Dự án Tốt nghiệp - Phân tích Hiệu suất và Dự báo Sản lượng Hệ thống Điện Mặt Trời},
    pdfauthor   = {The Outliers - FPT Polytechnic},
    bookmarksnumbered = true,
}

% --- CẤU HÌNH TIÊU ĐỀ CHƯƠNG ---
\usepackage{titlesec}
\titleformat*{\section}{\fontsize{13pt}{15.6pt}\bfseries}
\titleformat*{\subsection}{\fontsize{13pt}{15.6pt}\bfseries}
\titleformat*{\subsubsection}{\fontsize{13pt}{15.6pt}\bfseries}
\titleformat{\chapter}[display]
  {\normalfont\huge\bfseries}{}{0pt}{\huge}
\titlespacing*{\chapter}{0pt}{-10pt}{20pt} 

% ======================== HEADER / FOOTER ====================================
\fancypagestyle{main}{
    \fancyhf{}
    \fancyhead[L]{
\includegraphics[height=1cm]{figures/logo_fpoly.png}}
    \fancyhead[R]{\small THE OUTLIERS}
    \fancyfoot[L]{\small Dự án Tốt nghiệp — Chuyên ngành Xử lý Dữ liệu}
    \fancyfoot[R]{\small Trang \thepage\ / \pageref{LastPage}}
    \renewcommand{\headrulewidth}{0.4pt}
    \renewcommand{\footrulewidth}{0.4pt}
}

\fancypagestyle{plain}{
    \fancyhf{}
    \fancyfoot[L]{\small Dự án Tốt nghiệp — Chuyên ngành Xử lý Dữ liệu}
    \fancyfoot[R]{\small Trang \thepage\ / \pageref{LastPage}}
    \renewcommand{\headrulewidth}{0pt}
    \renewcommand{\footrulewidth}{0.4pt}
}

\fancypagestyle{front}{
    \fancyhf{}
    \fancyfoot[L]{\small Dự án Tốt nghiệp — Chuyên ngành Xử lý Dữ liệu}
    \fancyfoot[R]{\small Trang \thepage\ / \pageref{LastPage}}
    \renewcommand{\headrulewidth}{0pt}
    \renewcommand{\footrulewidth}{0.4pt}
}


\begin{document}
\pagestyle{front}

\begin{titlepage}
\enlargethispage{1.2cm}

\begin{tikzpicture}[remember picture, overlay]
    \draw[line width=2pt, fptorange]
        ([shift={(1cm,-1cm)}]current page.north west)
        rectangle
        ([shift={(-1cm,1cm)}]current page.south east);
    \draw[line width=0.5pt, fptorange]
        ([shift={(1.15cm,-1.15cm)}]current page.north west)
        rectangle
        ([shift={(-1.15cm,1.15cm)}]current page.south east);
\end{tikzpicture}

\begin{center}
    \textbf{\large TRƯỜNG CAO ĐẲNG FPT POLYTECHNIC}

    \vspace{0.5cm}
    
\includegraphics[width=0.38\textwidth]{figures/logo_fpoly.png}

    \vspace{0.4cm}
    {\Large \textbf{DỰ ÁN TỐT NGHIỆP}} \\[4pt]
    {\large \textbf{CHUYÊN NGÀNH: XỬ LÝ DỮ LIỆU}}

    \vspace{1.6cm}
    
    \begin{spacing}{1.0}
        {\huge \textbf{PHÂN TÍCH HIỆU SUẤT VÀ \\[4pt] DỰ BÁO SẢN LƯỢNG \\[4pt] HỆ THỐNG ĐIỆN MẶT TRỜI}}
    \end{spacing}

    \vspace{1.8cm}
    \normalsize
    \begin{flushleft}
        \hspace{2cm} \textbf{Nhóm thực hiện:} The Outliers \\[5pt]
        \hspace{2cm} \textbf{Sinh viên thực hiện:} \\[3pt]
        \hspace{2cm}
        \begin{tabular*}{7.5cm}{@{} l @{\extracolsep{\fill}} r @{}}
            \hspace{0.5cm} 1. Nguyễn Vũ Nhựt   & PS44442 \\
            \hspace{0.5cm} 2. Nguyễn Văn Sỹ    & PS43671 \\
            \hspace{0.5cm} 3. Lê Công Toàn    & PS44145 \\
            \hspace{0.5cm} 4. Nguyễn Xuân Hùng & PS44527 \\
            \hspace{0.5cm} 5. Ngô Tấn Đạt     & PS44589 \\
        \end{tabular*} \\[5pt]
        \hspace{2cm} \textbf{Giảng viên hướng dẫn:} Văn Công Khanh
    \end{flushleft}
    \vfill
    {\normalsize \textbf{TP. HỒ CHÍ MINH, THÁNG 6 NĂM 2026}}
    \vspace{0.5cm}
\end{center}
\end{titlepage}


\pagenumbering{arabic}
\pagestyle{main}

% Sinh tu Section4.tex - chi giu tu \chapter tro di.
\chapter{Phát hiện Dị thường và Làm sạch Dữ liệu Nâng cao}
\label{chap:outlier_cleaning}

\section{Tổng quan Kiến trúc Pipeline Xử lý Dữ liệu}
\label{sec:pipeline_overview}

Để xây dựng một hệ thống phân tích dữ liệu sản lượng điện mặt trời ổn định và tin cậy, dự án thiết lập quy trình trích xuất, biến đổi và nạp dữ liệu (ETL Pipeline) tự động hóa bằng ngôn ngữ Python. Luồng di chuyển dữ liệu (Data Flow) được chia làm 3 tầng lõi rõ rệt nhằm đáp ứng triết lý thiết kế kho dữ liệu tập trung (tiếp cận từ tổng quan đến chi tiết):

\begin{enumerate}[noitemsep]
    \item \textbf{Tầng Trích xuất (Extract):} Dữ liệu lịch sử từ các tệp tin CSV vật lý ghi nhận tại trạm pin và dữ liệu khí tượng viễn thám từ Open-Meteo API được thu thập và lưu trữ tập trung tại thùng chứa (Bucket) \texttt{raw-data} của hệ thống lưu trữ đối tượng MinIO S3. Các tệp tin cốt lõi bao gồm \texttt{campus\_meta.csv}, \texttt{Solar\_Site\_Details.csv}, \texttt{calender.csv}, \texttt{Solar\_Energy\_Generation.csv}, và \texttt{open\_meteo\_weather\_raw\_2020\_2022.csv}.
    \item \textbf{Tầng Biến đổi (Transform):} Sử dụng các thư viện tính toán hiệu năng cao như \texttt{pandas}, \texttt{numpy}, và \texttt{scikit-learn} để thực hiện chuẩn hóa kiểu dữ liệu thời gian, điền khuyết thiếu bằng chiến lược lai (điền khuyết kết hợp), gán cờ ngoại lai thông qua phân tích biến động trung bình trượt cục bộ, và đồng bộ granularity thời gian.
    \item \textbf{Tầng Nạp (Load):} Dữ liệu sau biến đổi được đẩy vào hệ quản trị cơ sở dữ liệu Supabase PostgreSQL thông qua hai schema chuyên biệt: \texttt{staging} (bộ đệm trung gian) và \texttt{datawarehouse} (kho dữ liệu chuẩn hóa đa chiều phục vụ BI và ML).
\end{enumerate}

Để khắc phục các hạn chế về môi trường triển khai thực tế, mã nguồn ETL của dự án tích hợp hai giải pháp hệ thống quan trọng:
\begin{itemize}
    \item \textbf{Sử dụng Thư viện kết nối thuần Python (\texttt{pg8000}):} Khác với thư viện \texttt{psycopg2} vốn yêu cầu biên dịch các thư viện hệ thống nhị phân C (dễ gây lỗi thiếu \texttt{libz.so.1} hoặc \texttt{libstdc++.so.6} trên môi trường ảo hóa NixOS), thư viện \texttt{pg8000} hoạt động hoàn toàn độc lập, đảm bảo tính di động cao của mã nguồn.
    \item \textbf{Cấu hình Connection Pooler thông qua IPv4:} Do Supabase chỉ hỗ trợ truy cập IPv6 trực tiếp cho gói tài khoản miễn phí, dự án cấu hình Pooler trung gian của Supabase tại cổng \texttt{5432} để duy trì kết nối IPv4 ổn định cho toàn bộ thành viên nhóm, đồng thời giảm tải luồng kết nối đồng thời lên máy chủ.
\end{itemize}

Dưới đây là mã nguồn cấu hình kết nối tối ưu hóa bằng thư viện \texttt{pg8000}:

\begin{lstlisting}[language=Python, caption={Mã nguồn thiết lập kết nối cơ sở dữ liệu qua Pooler bằng pg8000}, label={lst:pg8000_conn}]
def connect_pg8000() -> any:
    import pg8000.dbapi
    load_dotenv(ENV_FILE)
    required = [
        "DB_HOST", "DB_PORT", "DB_NAME", 
        "DB_USER", "DB_PASSWORD"
    ]
    missing = [name for name in required if not os.getenv(name)]
    if missing:
        missing_str = ', '.join(missing)
        err_msg = f"Missing vars in {ENV_FILE}: {missing_str}"
        raise RuntimeError(err_msg)

    return pg8000.dbapi.connect(
        host=os.environ["DB_HOST"],
        port=int(os.environ.get("DB_PORT", "5432")),
        database=os.environ["DB_NAME"],
        user=os.environ["DB_USER"],
        password=os.environ["DB_PASSWORD"],
        ssl_context=True,
    )
\end{lstlisting}

\subsection{Thiết lập Tiến trình Nạp Dữ liệu Thứ tự}
Để bảo toàn tính toàn vẹn tham chiếu khóa ngoại (Foreign Key Constraints), tiến trình tải dữ liệu vào kho được kiểm soát nghiêm ngặt theo thứ tự phụ thuộc. File \texttt{load\_01\_dims.py} chịu trách nhiệm tẩy trống dữ liệu cũ bằng lệnh \texttt{TRUNCATE CASCADE} và tải lại các bảng chiều độc lập theo trình tự tuyến tính:
\begin{equation*}
\text{dim\_geography} \longrightarrow \text{dim\_solar\_site} \longrightarrow \text{dim\_date} \longrightarrow \text{dim\_time} \longrightarrow \text{dim\_weather\_type}
\end{equation*}

Ngay sau khi các bảng chiều được nạp thành công, tệp tin \texttt{load\_02\_facts.py} tiến hành đọc các bảng sự kiện quy mô lớn từ MinIO và ánh xạ các khóa tự nhiên sang khóa đại diện (Surrogate Keys) bằng các bộ từ điển tra cứu (Lookups) đã lưu trong bộ nhớ tạm (In-memory dicts). Để tránh hiện tượng nghẽn bộ nhớ đệm (RAM Out-Of-Memory) khi xử lý tập dữ liệu sản lượng hơn 2,7 triệu dòng, luồng nạp sử dụng cơ chế chia nhỏ dữ liệu thành từng khối (chunksize = 50.000 dòng) và thực thi lệnh chèn hàng loạt \texttt{execute\_values} tối ưu hóa hiệu năng IO.

\subsection{Cơ chế Quản lý Giao dịch và An toàn Dữ liệu}
Toàn bộ tiến trình nạp được bọc trong một khối giao dịch cơ sở dữ liệu duy nhất (Explicit DB Transaction). Chế độ tự động xác nhận (\texttt{autocommit}) được tắt hoàn toàn. 
\begin{itemize}
    \item Nếu xuất hiện bất kỳ lỗi hệ thống hoặc vi phạm ràng buộc dữ liệu nào (ví dụ: khóa ngoại không tồn tại), lệnh \texttt{conn.rollback()} sẽ lập tức hoàn tác mọi thay đổi để đưa kho dữ liệu về trạng thái an toàn trước đó.
    \item Chỉ khi toàn bộ dữ liệu của phân đoạn được ghi thành công không lỗi, lệnh \texttt{conn.commit()} mới được phát ra để ghi nhận vĩnh viễn các thay đổi vào đĩa cứng.
\end{itemize}

\section{Đánh giá Chất lượng và Khám phá Dữ liệu (EDA ban đầu)}
\label{sec:data_quality}

Trước khi thực hiện các phép toán làm sạch, nhóm tiến hành một bước hồ sơ hóa dữ liệu (Data Profiling) nhằm phát hiện các vấn đề về chất lượng dữ liệu thô. Tiến trình này được tự động hóa thông qua kịch bản \texttt{data\_integrity\_and\_reconciliation\_check.py}. 

Công cụ thực hiện quét toàn bộ các bảng trong schema \texttt{staging} để thống kê: (1) Tổng số dòng dữ liệu hiện có, (2) Số lượng bản ghi bị khuyết thiếu khóa chính (NULL PK), và (3) Số lượng bản ghi bị trùng lặp khóa tự nhiên (Duplicate PKs). 

\begin{table}[H]
\centering
\caption{Bảng thống kê kiểm tra chất lượng dữ liệu ban đầu trong Staging}
\label{tab:data_quality_report}
\renewcommand{\arraystretch}{1.3}
% Sử dụng >{\centering\arraybackslash}X để biến các cột số liệu thành cột tự động căn giữa và tự động xuống dòng
\begin{tabularx}{\textwidth}{l >{\centering\arraybackslash}X >{\centering\arraybackslash}X >{\centering\arraybackslash}X}
\toprule
\textbf{Tên bảng dữ liệu} & \textbf{Tổng số dòng} & \textbf{Số dòng trùng lặp} & \textbf{Số dòng khuyết PK} \\
\midrule
\texttt{dim\_geography}    & 5          & 0                  & 0                 \\
\texttt{dim\_solar\_site}  & 42         & 0                  & 0                 \\
\texttt{dim\_date}         & 2.312       & 0                  & 0                 \\
\texttt{dim\_time}         & 96         & 0                  & 0                 \\
\texttt{dim\_weather\_type}& 22         & 0                  & 0                 \\
\texttt{fact\_weather}     & 367.920     & 0                  & 0                 \\
\texttt{fact\_solar\_energy\_gen} & 2.731.946 & 0            & 0                 \\
\bottomrule
\end{tabularx}
\end{table}

Kết quả phân tích ban đầu cho thấy cấu trúc khung của các bảng chiều và bảng thời tiết rất sạch, không xuất hiện hiện tượng trùng lặp dòng. Tuy nhiên, bảng sự kiện sản lượng thô \texttt{fact\_solar\_energy\_gen} chứa tới 107.452 ô dữ liệu bị khuyết giá trị sản lượng (\texttt{SolarGeneration} mang giá trị \texttt{NaN}), chiếm khoảng 3,93\% tổng số quan trắc đo đạc của cả năm. Đây là đối tượng trọng tâm cần xử lý tại tầng làm sạch dữ liệu.

\subsection{Các Cột Thực Tế và Ánh Xạ Dữ Liệu Trong Pipeline}
Trong giai đoạn nạp dữ liệu từ các tệp thô vào bảng đệm (Staging layer), hệ thống thực hiện chuyển đổi kiểu dữ liệu (casting) và đổi tên trường để phù hợp với định dạng kho dữ liệu. Chi tiết các trường thông tin thực tế được ánh xạ từ tệp thô lên bảng sự kiện được thể hiện qua mã nguồn xử lý trực tiếp dưới đây:

\subsubsection{Cấu trúc ánh xạ bảng sự kiện sản lượng (\texttt{fact\_solar\_energy\_gen})}
Bảng này ghi nhận sản lượng điện thực tế theo chu kỳ 15 phút của từng trạm phát:
\begin{itemize}
    \item \texttt{sitekey} (Ánh xạ từ cột \texttt{SiteKey} trong tệp thô, định dạng chuỗi chuyển sang \texttt{VARCHAR} hoặc \texttt{INT}): Mã định danh duy nhất của từng trạm phát.
    \item \texttt{timestamp} (Ánh xạ từ cột \texttt{Timestamp} dạng text thô, xử lý chuyển đổi định dạng \texttt{YYYY-MM-DD HH24:MI:SS} sang kiểu dữ liệu \texttt{TIMESTAMP} chuẩn): Mốc thời gian quan trắc sản lượng.
    \item \texttt{energy\_generated\_kwh} (Ánh xạ từ cột \texttt{SolarGeneration} thô, làm sạch ký tự lạ và chuyển đổi thành \texttt{DOUBLE PRECISION} hoặc \texttt{FLOAT}): Sản lượng điện mặt trời thực tế thu được trong 15 phút (đơn vị kWh).
\end{itemize}

\subsubsection{Cấu trúc ánh xạ bảng sự kiện thời tiết (\texttt{fact\_weather})}
Thay vì chỉ liệt kê thủ công, toàn bộ logic trích xuất, ép kiểu dữ liệu từ nguồn khí tượng thô \texttt{staging.stg\_open\_meteo\_weather\_raw} (lấy từ Open-Meteo API~\cite{openmeteo2023}) sang bảng sự kiện \texttt{fact\_weather} được hiện thực hóa qua định nghĩa \texttt{LoadSpec} trong tệp tin mã nguồn \texttt{srcs/02\_transform/01\_run\_transform\_buffers.py}:

\begin{lstlisting}[language=Python, caption={Mã nguồn Python/SQL thực tế ánh xạ và chuyển đổi kiểu dữ liệu cho bảng sự kiện fact\_weather}, label={lst:fact_weather_mapping}]
LoadSpec(
    table="fact_weather",
    columns=(
        "sitekey",
        "timestamp",
        "weather_code",
        "is_day",
        "shortwave_radiation",
        "temperature_c",
        "cloud_cover_total",
        "cloud_cover_low",
        "cloud_cover_mid",
        "cloud_cover_high",
        "diffuse_solar_radiation",
        "direct_normal_irradiance",
        "wind_speed",
        "precipitation_mm",
        "sunshine_duration",
    ),
    select_sql=\"\"\"
        SELECT
            sitekey,
            timestamp::timestamp AS timestamp,
            weather_code::int AS weather_code,
            is_day::int AS is_day,
            shortwave_radiation::numeric 
                AS shortwave_radiation,
            temperature_2m::numeric 
                AS temperature_c,
            cloud_cover::numeric 
                AS cloud_cover_total,
            cloud_cover_low::numeric 
                AS cloud_cover_low,
            cloud_cover_mid::numeric 
                AS cloud_cover_mid,
            cloud_cover_high::numeric 
                AS cloud_cover_high,
            diffuse_radiation::numeric 
                AS diffuse_solar_radiation,
            direct_radiation::numeric 
                AS direct_normal_irradiance,
            wind_speed_10m::numeric 
                AS wind_speed,
            precipitation::numeric 
                AS precipitation_mm,
            sunshine_duration::numeric 
                AS sunshine_duration
        FROM staging.stg_open_meteo_weather_raw
    \"\"\",
)
\end{lstlisting}

Đoạn mã trên thể hiện rõ quy tắc chuẩn hóa kiểu dữ liệu nghiêm ngặt: mốc thời gian được ép về \texttt{TIMESTAMP}, mã thời tiết \texttt{weather\_code} và cờ \texttt{is\_day} ép về \texttt{INT}, trong khi tất cả 11 biến khí tượng viễn thám liên tục (bức xạ mặt trời, nhiệt độ 2m, các tầng mây che phủ, tốc độ gió 10m, lượng mưa và thời lượng nắng) đều được ép kiểu \texttt{NUMERIC} để bảo toàn độ chính xác số học cao nhất cho các phép tính tổn thất và mô hình học máy.

\section{Làm sạch và Chuẩn hóa Dữ liệu (Data Cleaning)}
\label{sec:cleaning}

\subsection{Xử lý Dữ liệu Khuyết thiếu và Lỗi Logic}
\label{subsec:null_imputation}

Để phục hồi dữ liệu khuyết thiếu cho trường sản lượng phát điện (\texttt{energy\_generated\_kwh}), dự án áp dụng chiến lược nội suy lai nghiêm ngặt (Hybrid Imputation Strategy) tích hợp giữa các bộ quy tắc nghiệp vụ vật lý và mô hình học máy trong file \texttt{solar\_data\_pipeline.py}. Tiến trình xử lý gồm hai giai đoạn chính:

\subsubsection{Giai đoạn 1: Lọc vật lý ban đêm (Rule-Based Night Zeroing)}
Đối với các trạm phát điện mặt trời quang điện, sản lượng điện năng luôn bằng 0 vào ban đêm do không có bức xạ mặt trời. Nếu sử dụng các thuật toán nội suy thuần toán học (như tuyến tính hay spline) trên các khoảng trống kéo dài qua đêm, hệ thống sẽ sinh ra các giá trị sản lượng ảo (dương) phi vật lý. 

Do đó, thuật toán tiến hành rà soát mốc thời gian và áp đặt giá trị 0 ngay lập tức nếu thỏa mãn một trong hai điều kiện logic sau:
\begin{itemize}
    \item \textbf{Khung giờ ban đêm tự nhiên:} Mốc thời gian đo đạc nằm trong khoảng từ 18:30 chiều hôm trước đến 05:30 sáng hôm sau ($T \ge 18.5$ hoặc $T < 5.5$).
    \item \textbf{Điều kiện thời tiết triệt tiêu:} Giá trị bức xạ sóng ngắn mặt trời (\texttt{shortwave\_radiation}) bằng 0 hoặc cờ chỉ số ánh sáng ngày \texttt{is\_day} bằng 0 trong tập dữ liệu thời tiết (được đồng bộ hóa sang dòng sản lượng tương ứng bằng thuật toán ghép nối mốc gần nhất \texttt{pd.merge\_asof} với ngưỡng sai lệch thời gian tối đa là 30 phút).
\end{itemize}
Giai đoạn này đã xử lý thành công 94.618 ô dữ liệu trống (chiếm tới 88,1\% tổng số giá trị khuyết thiếu ban đầu), gán chúng về giá trị \texttt{0.0} một cách chính xác.

Mã nguồn Python hiện thực thuật toán lọc vật lý ban đêm:

\begin{lstlisting}[language=Python, caption={Mã nguồn Python hiện thực bộ quy tắc Night-Zeroing vật lý}, label={lst:night_zero}]
def rule_based_night_zero(
    solar_df: pd.DataFrame, weather_df: pd.DataFrame, site_key: int
) -> tuple:
    gen = solar_df["energy_generated_kwh"].copy()
    ts = solar_df["timestamp"]

    # Xac dinh gio dem tu nhien
    hour = ts.dt.hour + ts.dt.minute / 60
    is_night = (hour >= NIGHT_START + 0.5) | (hour < NIGHT_END + 0.5)

    # Ket noi voi open_meteo de lay thong tin buc xa song ngan
    w = weather_df[weather_df["sitekey"] == site_key][
        ["timestamp", "shortwave_radiation", "is_day"]
    ].copy()
    w = w.rename(columns={"timestamp": "timestamp"})
    w["timestamp"] = pd.to_datetime(w["timestamp"])

    merged = pd.merge_asof(
        solar_df[["timestamp"]].reset_index(),
        w.sort_values("timestamp"),
        on="timestamp",
        tolerance=pd.Timedelta("30min"),
        direction="nearest",
    ).set_index("index")

    # Dieu kien buc xa song ngan bang 0 hoac co dem is_day = 0
    zero_radiation = (merged["shortwave_radiation"].fillna(0) == 0) | (
        merged["is_day"].fillna(0) == 0
    )
    mask_zero = is_night.values | zero_radiation.values

    n_filled = (mask_zero & gen.isna()).sum()
    gen[mask_zero & gen.isna()] = 0.0
    return gen, n_filled
\end{lstlisting}

\subsubsection{Giai đoạn 2: Nội suy phân cấp theo quy mô khoảng trống (Imputation Gaps classification)}
Đối với 12.834 ô dữ liệu khuyết thiếu phát sinh vào ban ngày (khi có bức xạ), thuật toán tiến hành phân tích độ dài chuỗi khuyết liên tiếp (Gap size) của từng trạm phát độc lập để áp dụng phương pháp điền khuyết tối ưu:

\begin{enumerate}[label=\alph*., leftmargin=*]
    \item \textbf{Khoảng khuyết ngắn ($\le 2$ dòng liên tiếp, tương đương $\le 30$ phút):} Áp dụng phương pháp nội suy tuyến tính (\textit{Linear Interpolation}) theo thời gian. Do khoảng thời gian gián đoạn rất ngắn, biến động của ánh sáng mặt trời được coi là thay đổi tuyến tính đều giữa điểm biên trước và biên sau khoảng khuyết.
    \item \textbf{Khoảng khuyết trung bình (Từ 3 đến 8 dòng liên tiếp, tương đương từ 45 phút đến 2 tiếng):} Áp dụng phương pháp nội suy \textit{Cubic Spline}. Đường cong bậc ba giúp mô phỏng mượt mà sự trồi sụt tự nhiên của sản lượng điện do mây che phủ đi qua trạm pin. Để tránh hiện tượng sập nhiễu ma trận khi gặp các khoảng khuyết lớn lân cận, thuật toán xây dựng một bản sao tạm thời được điền tuyến tính để làm mịn dữ liệu mốc trước khi tính toán spline bậc ba thực tế.
    \item \textbf{Khoảng khuyết diện rộng ($> 8$ dòng liên tiếp, tương đương $> 2$ tiếng):} Trong trường hợp này, dữ liệu bị mất quá dài khiến các phương pháp nội suy chuỗi đơn biến (Tuyến tính/Spline) mất khả năng phục hồi chính xác xu hướng thực tế. Nhóm thiết lập mô hình hồi quy đa biến (\textit{Multivariate Regression Imputation}) độc lập cho từng trạm phát.
    
    Mô hình Hồi quy Tuyến tính (\texttt{LinearRegression} từ thư viện \texttt{scikit-learn}) được huấn luyện trên tập các bản ghi sạch (đầy đủ thông tin) của chính trạm đó. Các biến đầu vào (Features) bao gồm: bức xạ sóng ngắn mặt trời (\texttt{shortwave\_radiation}), bức xạ chiếu trực tiếp (\texttt{direct\_normal\_irradiance}), bức xạ khuếch tán (\texttt{diffuse\_solar\_radiation}) và nhiệt độ không khí ở độ cao 2m (\texttt{temperature\_c}). 
    
    Sau khi huấn luyện, mô hình sẽ dự đoán sản lượng điện năng cho các khoảng khuyết diện rộng dựa trên dữ liệu khí tượng thực tế tại tọa độ trạm phát đó. Kết quả dự đoán được cắt giới hạn dưới bằng 0 (\texttt{clip(lower=0)}) để loại bỏ các giá trị âm sinh ra do sai số hồi quy.
\end{enumerate}

Mã nguồn Python thực thi hồi quy đa biến điền khuyết diện rộng ban ngày:

\begin{lstlisting}[language=Python, caption={Mã nguồn Python thực thi Multivariate Regression Imputation cho các khoảng khuyết lớn}, label={lst:regression_impute}]
def regression_imputation_large_gaps_strict(
    solar_sub: pd.DataFrame,
    weather_df: pd.DataFrame,
    site_key: int,
    large_gap_indices: set,
) -> tuple:
    FEATURES = [
        "shortwave_radiation",
        "direct_normal_irradiance",
        "diffuse_solar_radiation",
        "temperature_c",
    ]

    cols = ["timestamp"] + FEATURES
    w = weather_df[weather_df["sitekey"] == site_key][cols].copy()
    w = w.rename(columns={"timestamp": "timestamp"})
    w["timestamp"] = pd.to_datetime(w["timestamp"])

    merged = pd.merge_asof(
        solar_sub[["timestamp", "energy_generated_kwh"]].reset_index(),
        w.sort_values("timestamp"),
        on="timestamp",
        tolerance=pd.Timedelta("60min"),
        direction="nearest",
    ).set_index("index")

    train_mask = merged["energy_generated_kwh"].notna() & \
                 merged[FEATURES].notna().all(axis=1)
    pred_mask = merged["energy_generated_kwh"].isna() & \
                merged[FEATURES].notna().all(axis=1)

    filled = 0
    if train_mask.sum() < 10:
        return solar_sub["energy_generated_kwh"].values, 0

    X_train = merged.loc[train_mask, FEATURES].values
    y_train = merged.loc[train_mask, "energy_generated_kwh"].values

    # Chuan hoa thong so dac trung dau vao
    scaler = StandardScaler()
    X_train_sc = scaler.fit_transform(X_train)

    model = LinearRegression()
    model.fit(X_train_sc, y_train)

    if pred_mask.sum() > 0:
        X_pred = merged.loc[pred_mask, FEATURES].values
        X_pred_sc = scaler.transform(X_pred)
        y_pred = model.predict(X_pred_sc)

        result = solar_sub["energy_generated_kwh"].copy()
        pred_indices = merged[pred_mask].index.tolist()

        for pos, idx in enumerate(pred_indices):
            local_pos = (
                solar_sub.index.get_loc(idx)
                if idx in solar_sub.index else None
            )
            # Chi ghi vao cac vi tri thuoc khoang khuyet lon
            if local_pos is not None and local_pos in large_gap_indices:
                result.iloc[local_pos] = max(0.0, y_pred[pos])
                filled += 1

        return result.values, filled
    return solar_sub["energy_generated_kwh"].values, 0
\end{lstlisting}

Tổng hợp kết quả sau quy trình làm sạch dữ liệu khuyết thiếu:
\begin{itemize}
    \item Tổng số ô trống ban đầu: 107.452 ô (3,93\% dữ liệu).
    \item Số ô được điền qua lọc ban đêm (Night Zero): 94.618 ô.
    \item Số ô được điền qua nội suy tuyến tính (Linear): 3.652 ô.
    \item Số ô được điền qua nội suy đường cong (Cubic Spline): 5.124 ô.
    \item Số ô được điền qua mô hình học máy hồi quy (Regression): 4.058 ô.
    \item Số ô NULL còn lại trong kho dữ liệu: 0 ô (Điền khuyết thành công 100\%).
\end{itemize}

\subsection{Xử lý Dữ liệu Ngoại lai (Outliers) và Lọc nhiễu}
\label{subsec:outlier_processing}

Trong chuỗi thời gian vận hành thực tế của hệ thống điện mặt trời quang điện (PV), dữ liệu đo đạc từ các cảm biến dòng điện (CT) và đồng hồ đo năng lượng (Smart Meters) thường xuyên xuất hiện các điểm dị thường (Anomalies/Outliers). Nếu không được nhận diện và gắn cờ chính xác, các dữ liệu này sẽ làm biến dạng các chỉ số hiệu suất kỹ thuật (như PR, CF) và gây sai lệch nghiêm trọng cho mô hình huấn luyện học máy dự báo sản lượng.

\subsubsection{Bản chất Nghiệp vụ và Phân tích Nguyên nhân Phát sinh Ngoại lai}
Dựa trên các nghiên cứu kỹ thuật chuyên sâu về độ tin cậy quang điện áp mái từ Tổ chức Năng lượng Quốc tế (IEA-PVPS Task 13)~\cite{iea2023pvps} và khảo sát thực nghiệm tại 42 trạm PV của Đại học La Trobe (Úc)~\cite{wimalaratne2022unisolar}, các hiện tượng bất thường trong dữ liệu phát sinh từ 4 nhóm nguyên nhân vật lý cốt lõi:

\begin{enumerate}[label=\arabic*., leftmargin=*]
    \item \textbf{Sự cố Biến tần và Thiết bị Chuyển đổi (Inverter / Optimizer Failures):}
    Theo báo cáo của IEA-PVPS~\cite{iea2023pvps}, bộ biến tần (Inverter) là mắt xích có độ tin cậy thấp nhất trong toàn bộ vòng đời hệ thống PV. Thời gian trung bình giữa hai lần hỏng (MTBF) của Inverter ngắn hơn tấm pin quang điện từ 300 đến 500 lần. Lỗi Inverter chiếm tới \textbf{36\%} tổng sản lượng tổn thất do thiết bị vật lý gây ra, với tỷ lệ hư hỏng hàng năm dao động từ 1\% đến 15\%/năm. Khi Inverter bị ngắt mạch bảo vệ quá nhiệt (Thermal trip) hoặc hỏng linh kiện bán dẫn, sản lượng trạm sụt giảm đột ngột về 0 giữa ban ngày dù bức xạ mặt trời đang đạt đỉnh cực đại ($> 700 \, \text{W/m}^2$).

    \item \textbf{Hỏng Diode Bypass và Che bóng Cục bộ (Bypass Diode Faults \& Shading):}
    Khi một tấm pin trong chuỗi (String) bị che bóng bởi cây cối, bụi bẩn bám dính dày đặc (Soiling) hoặc hỏng diode bypass, dòng điện của toàn bộ chuỗi bị bóp nghẹt theo tấm pin yếu nhất. Sự cố diode bypass có thể làm triệt tiêu tới \textbf{33\%} sản lượng điện của riêng tấm pin đó, tạo nên các vết sụt giảm sản lượng cục bộ dạng bậc thang phi tự nhiên.

    \item \textbf{Suy giảm do Quá nhiệt và Điểm nóng (Thermal Derating \& Hot-spots):}
    Hệ số nhiệt độ (Temperature Coefficient) của pin silicon tinh thể dao động từ $-0{,}30\%$ đến $-0{,}50\%/^\circ\text{C}$ khi nhiệt độ vượt mốc điều kiện chuẩn STC ($25^\circ\text{C}$). Vào các ngày nắng gắt mùa hè, nhiệt độ bề mặt ô pin ($T_{\text{cell}}$) có thể tăng vọt lên tới $65^\circ\text{C}$, làm sụt giảm \textbf{12--20\%} công suất đỉnh chủ yếu do suy giảm điện áp hở mạch ($V_{oc}$). Nhiệt độ vượt ngưỡng $85--90^\circ\text{C}$ có thể gây ra hiện tượng điểm nóng (Hot-spots), dẫn đến suy thoái vĩnh viễn mối hàn và lớp phủ bảo vệ.

    \item \textbf{Nhiễu Cảm biến và Dòng điện Rò Ban đêm (Sensor Drift \& Dòng rò ban đêm):}
    Hiện tượng rò rỉ dòng điện ngược (Reverse Current) từ lưới điện vào ban đêm hoặc cảm biến biến dòng (CT) bị trôi mốc 0 (Zero-drift fault) sinh ra các giá trị sản lượng dương giả mạo ($E > 0$) trong khoảng thời gian không có bức xạ ($18:30 - 05:30$). Ngoài ra, các xung nhiễu điện áp lưới (Voltage spikes) hoặc gián đoạn đường truyền viễn thông (Telemetry dropouts) tạo nên các đỉnh xung sản lượng vượt trần công suất thiết kế cực đại ($P > 1{,}2 \cdot P_{\text{stc}}$) hoặc các bước nhảy đột biến giữa các chu kỳ 15 phút kế cận.
\end{enumerate}

\subsubsection{Khung Thuật toán Phân lớp Lai GMM--IF và Rào chắn Vật lý}
Dữ liệu chuỗi thời gian quang điện có tính phi tuyến và phân phối đa đỉnh (Multimodal) biến động theo thời tiết, khiến các phương pháp thống kê toàn cục kinh điển (như $3\sigma$ hay Boxplot IQR truyền thống~\cite{tukey1977eda}) dễ nhận diện nhầm dữ liệu bình thường của ngày âm u thành ngoại lai. Nhằm khắc phục hạn chế này, nhóm mở rộng khung nghiên cứu của Li và cộng sự~\cite{li2025gmmif} trên tạp chí \textit{IEEE Access}, thiết lập mô hình lai phân lớp 3 tầng trong \texttt{srcs/02\_transform/02\_generate\_outliers/02\_gmm\_if.py}:

\begin{itemize}
    \item \textbf{Tầng 1 -- Phân đoạn Cây quyết định (Decision Tree Segmentation):} Sử dụng cây hồi quy quyết định để chia dữ liệu của từng trạm phát thành 19 đến 29 vùng vận hành lá (Leaf segments) dựa trên bức xạ sóng ngắn, nhiệt độ và mốc thời gian. Mỗi phân đoạn tạo thành một vùng dữ liệu cục bộ tương đối đồng nhất có phân phối xấp xỉ Gauss (hệ số $R^2$ trung bình đạt $0{,}758$).
    \item \textbf{Tầng 2 -- Ước lượng Mật độ Hỗn hợp Gauss (GMM)~\cite{bishop2006pattern}:} Trong từng phân đoạn lá, huấn luyện mô hình GMM 2 thành phần ($M = 2$) để đo lường mật độ xác suất $p(x)$. Các điểm dữ liệu nằm ở vùng mật độ cực thấp ($P_{\text{GMM}} < 0{,}02$) bị gắn nhãn ứng viên ngoại lai theo góc nhìn phân phối xác suất.
    \item \textbf{Tầng 3 -- Rừng cô lập (Isolation Forest -- IF)~\cite{liu2008isolation}:} Khởi tạo 100 cây ngẫu nhiên để đo lường số nhát cắt cần thiết nhằm cô lập điểm dữ liệu khỏi đám đông với tham số \texttt{contamination} $= 3\%$.
    \item \textbf{Cơ chế Hợp nhất Đồng thuận (Fusion Consensus $GMM \wedge IF$):} Một điểm dữ liệu chỉ được gắn cờ dị thường máy học khi \textbf{cả GMM và IF cùng đồng thuận}:
    \begin{equation}
        \text{Flag}_{\text{ML}} = \text{Flag}_{\text{GMM}} \wedge \text{Flag}_{\text{IF}}
    \end{equation}
    Phép giao logic này giúp triệt tiêu phần lớn các cảnh báo giả (False Positives) của từng mô hình riêng lẻ (chỉ số tương đồng Jaccard giữa hai thuật toán đạt từ $6{,}4\%$ đến $16{,}7\%$).
\end{itemize}

\begin{figure}[H]
\centering
\includegraphics[width=0.98\textwidth]{figures/gmm_if_architecture.png}
\caption{Sơ đồ kiến trúc mô hình phân lớp lai GMM--IF và rào chắn vật lý kết hợp quy trình 4 bước nạp cờ an toàn vào Supabase DWH.}
\label{fig:gmm_if_architecture}
\end{figure}

\subsubsection{Hệ thống Danh mục Mã Lý do Ngoại lai (\texttt{gmm\_if\_outlier\_reason})}
Để phục vụ cho hệ thống giám sát và bảo trì cảnh báo trên Dashboard Tableau, mỗi bản ghi bị gắn cờ được gán một mã nguyên nhân cụ thể thông qua hàm \texttt{explain\_outlier\_row()} trong \texttt{02\_gmm\_if.py}. Chi tiết 6 mã lý do kỹ thuật được chuẩn hóa tại Bảng~\ref{tab:outlier_reasons_dict}:

\begin{table}[H]
\centering
\small
\renewcommand{\arraystretch}{1.3}
\begin{tabularx}{\textwidth}{@{} l p{4.5cm} X @{}}
\toprule
\textbf{Mã Định danh Lý do} & \textbf{Định nghĩa Ngưỡng Kỹ thuật} & \textbf{Ý nghĩa Nghiệp vụ \& Chẩn đoán} \\
\midrule
\texttt{GMM\_IF\_CONSENSUS} & 
$\text{Flag}_{\text{GMM}} = \text{true}$ \textbf{và} $\text{Flag}_{\text{IF}} = \text{true}$ (Phép giao đồng thuận). & 
Bất thường đồng thời về mặt mật độ xác suất cục bộ và không gian phân tách toàn cục; cần kiểm tra cảm biến hoặc trạng thái vận hành trạm. \\
\midrule
\texttt{PHYSICAL\_OVER\_CAPACITY} & 
Sản lượng 15 phút vượt trần công suất lắp đặt: $E > \text{capacity\_kw} \times 0{,}25$. & 
Đỉnh phát điện vượt quá công suất định mức cực đại của dàn PV; cảnh báo xung điện lưới, lỗi telemetry hoặc sai lệch siêu dữ liệu trạm. \\
\midrule
\texttt{PHYSICAL\_HIGH\_ENERGY\_NO\_SUN} & 
$\text{shortwave\_radiation} \le 25 \, \text{W/m}^2$, $\text{sunshine} \le 60\text{s}$, nhưng $E \ge \max(1.0, 0{,}20 \cdot P_{\text{stc}})$. & 
Sản lượng điện cao trong điều kiện trời tối hoàn toàn; mâu thuẫn vật lý cảnh báo sự cố rò rỉ điện ban đêm (Dòng rò ban đêm) hoặc dòng rò ngược. \\
\midrule
\texttt{PHYSICAL\_HIGH\_ENERGY\_LOW\_RAD} & 
$\text{shortwave\_radiation} \le 50 \, \text{W/m}^2$, $E \ge \max(1.0, 0{,}20 \cdot P_{\text{stc}})$, và $E > Q_3 + 4 \cdot \text{IQR}$. & 
Sản lượng cao bất thường khi bức xạ mặt trời rất yếu; vi phạm ngưỡng phân phối thống kê đuôi sâu của nhóm bức xạ tương ứng. \\
\midrule
\texttt{PHYSICAL\_LOW\_ENERGY\_STRONG\_SUN} & 
$\text{shortwave\_radiation} \ge 700 \, \text{W/m}^2$, $\text{sunshine} \ge 3000\text{s}$, $E \le 0{,}05 \cdot P_{95}$ và $E \le Q_1 - 2 \cdot \text{IQR}$. & 
Mất phát hoặc sụt giảm sản lượng nghiêm trọng trong điều kiện nắng cực đại; cảnh báo sự cố ngắt Inverter, hỏng diode bypass hoặc che bóng diện rộng. \\
\midrule
\texttt{PHYSICAL\_DISTRIBUTION\_JUMP} & 
Bỏ qua giờ chuyển tiếp (05h, 06h, 18h). $|E - E_{\text{neighbor}}| \ge \max(0{,}15 \cdot P_{95}, 1.0)$ và $E \notin [Q_1 - 4\text{IQR}, Q_3 + 4\text{IQR}]$. & 
Bước nhảy hoặc sụt giảm công suất cục bộ đột ngột so với diễn biến 2 giờ xung quanh; cảnh báo hiện tượng spike tín hiệu hoặc dropout truyền thông tin. \\
\bottomrule
\end{tabularx}
\caption{Danh mục các mã lý do ngoại lai được định nghĩa trong pipeline GMM--IF và ý nghĩa chẩn đoán bảo trì.}
\label{tab:outlier_reasons_dict}
\end{table}

Khi một bản ghi vi phạm đồng thời nhiều điều kiện, các mã lý do được nối với nhau bằng dấu cộng (Ví dụ: \texttt{GMM\_IF\_CONSENSUS+PHYSICAL\_DISTRIBUTION\_JUMP}), hỗ trợ kỹ sư vận hành phân rã đa tầng nguyên nhân sự cố.

\subsubsection{Quy trình Đẩy Cờ An toàn 4 Bước vào Kho Dữ liệu Supabase}
Để cập nhật cờ ngoại lai lên Supabase DWH mà không làm gián đoạn hệ thống và đảm bảo an toàn giao dịch tuyệt đối, nhóm xây dựng pipeline 4 bước trong mã nguồn \texttt{srcs/02\_transform/02\_run\_apply\_outlier\_flags.py}:
\begin{enumerate}
    \item \textbf{Xác thực nguồn dữ liệu (Source Verification Check):} Trước khi ghi dữ liệu, hàm \texttt{fingerprint\_supabase()} tiến hành đối chiếu cấu trúc vân tay dữ liệu giữa tệp kết quả phân tích ngoại lai và bảng sự kiện \texttt{fact\_solar\_energy\_gen} trên Supabase. 
    
    Quy trình kiểm tra tính toán tổng băm MD5 của chuỗi khóa ghép (\texttt{sitekey|timestamp}), tổng sản lượng năng lượng (\texttt{energy\_sum}) với sai số cho phép $\le 0{,}01$, và tổng số dòng dữ liệu. Bước này đảm bảo dữ liệu chạy cờ ngoại lai và dữ liệu trên cơ sở dữ liệu hoàn toàn đồng nhất về cấu trúc cột và số lượng dòng (2.731.946 dòng), ngăn ngừa lỗi lệch dòng (Data Shift).
    
    \item \textbf{Nạp danh sách ứng viên (Sparse Candidate Upload):} Thay vì tải lên lại toàn bộ 2,73 triệu dòng dữ liệu (rất chậm và dễ gây nghẽn mạng), thuật toán chỉ lọc ra 33.280 bản ghi bị gắn cờ ngoại lai thực sự (\texttt{gmm\_if\_outlier\_flag = true}). Danh sách ứng viên này được đẩy vào bảng tạm trung gian: \texttt{staging.temp\_gmm\_if\_outlier\_flags}.
    
    \item \textbf{Kiểm tra tính toàn vẹn khóa ngoại (Referential Check):} Thực hiện truy vấn kiểm tra chéo (Anti-Join) để đảm bảo toàn bộ các khóa trong bảng tạm đều tồn tại khớp 1:1 trong bảng sự kiện chính. Kết quả trả về 0 dòng lỗi trước khi tiến hành cập nhật.
    
    \item \textbf{Cập nhật giao dịch hàng loạt (Bulk Join Update):} Tiến trình mở một khối giao dịch an toàn trên cơ sở dữ liệu, thực hiện gán toàn bộ cờ \texttt{gmm\_if\_outlier\_flag = false} cho cả bảng sự kiện chính (2.731.946 dòng), sau đó thực hiện lệnh cập nhật \texttt{UPDATE ... FROM ... WHERE} nối bảng tạm với bảng chính theo khóa trùng khớp (\texttt{sitekey}, \texttt{timestamp}) để chuyển 33.280 dòng ngoại lai sang trạng thái \texttt{true} kèm mã lý do tương ứng.
\end{enumerate}

Dưới đây là một phần mã nguồn Python thực hiện so sánh đối chiếu vân tay (Fingerprint) dữ liệu trích từ \texttt{srcs/02\_transform/02\_run\_apply\_outlier\_flags.py}:

\begin{lstlisting}[language=Python, caption={Mã nguồn Python kiểm tra vân tay dữ liệu (Fingerprint) trước khi cập nhật}, label={lst:data_fingerprint}]
def fingerprint_supabase(conn: any) -> dict[str, object]:
    cur = conn.cursor()
    try:
        cur.execute(
            f"""
            SELECT
                COUNT(*)::bigint AS rows,
                COUNT(DISTINCT sitekey)::bigint AS sites,
                MIN(to_char(
                    timestamp, 'YYYY-MM-DD HH24:MI:SS'
                )),
                MAX(to_char(
                    timestamp, 'YYYY-MM-DD HH24:MI:SS'
                )),
                SUM(
                    energy_generated_kwh::numeric
                ) AS energy_sum,
                SUM(
                    ('x' || substr(
                        md5(sitekey::text || '|' || 
                        to_char(
                            timestamp, 
                            'YYYY-MM-DD HH24:MI:SS'
                        )),
                        1, 15
                    ))::bit(60)::bigint
                )::numeric AS key_checksum
            FROM {qname(SCHEMA, SOURCE_FACT_TABLE)}
            """
        )
        row = cur.fetchone()
        return {
            "rows": int(row[0]),
            "distinct_sitekeys": int(row[1]),
            "min_timestamp": row[2],
            "max_timestamp": row[3],
            "energy_sum": Decimal(str(row[4] or 0)),
            "key_checksum": Decimal(str(row[5] or 0)),
        }
    finally:
        cur.close()
\end{lstlisting}

\subsubsection{Kết quả Phân bổ và Đánh giá Kiểm toán Không Giám sát}
Vì đây là bài toán học máy không giám sát trên dữ liệu chuỗi thời gian thực tế không có nhãn sự cố thủ công từ người vận hành, hệ thống không thể sử dụng các chỉ số giả định như Accuracy, Precision hay Recall. Thay vào đó, chất lượng của bộ phát hiện ngoại lai được kiểm chứng thông qua bộ 4 chỉ số kiểm toán khoa học:

\begin{itemize}
    \item \textbf{Độ hội tụ mô hình GMM:} Đạt \textbf{100\%} trên toàn bộ 42 trạm phát (toàn bộ các phân đoạn lá cây quyết định đều hội tụ thuật toán EM thành công, không có segment nào bị lỗi fit).
    \item \textbf{Kiểm soát tỷ lệ nhiễu IF:} Isolation Forest giữ đúng tham số \texttt{contamination} xấp xỉ \textbf{3\%} trên từng trạm phát.
    \item \textbf{Chất lượng phân đoạn cây quyết định ($R^2$):} Hệ số $R^2$ đạt trung bình $\mathbf{0{,}758}$ (thấp nhất $0{,}606$, cao nhất $0{,}826$), chứng minh các vùng phân đoạn mô phỏng chính xác cấu trúc hình dạng phát điện của trạm.
    \item \textbf{Tổng số lượng bản ghi gắn cờ ngoại lai:} Toàn bộ pipeline đã gắn cờ thành công $\mathbf{33.280}$ dòng trên tổng số $\mathbf{2.731.946}$ dòng dữ liệu, tương ứng với tỷ lệ ngoại lai được kiểm soát ở mức lý tưởng $\mathbf{1{,}22\%}$.
\end{itemize}

\section{Biến đổi và Trích xuất Đặc trưng (trích xuất đặc trưng nghiệp vụ)}
\label{sec:transformation}

Sau khi dữ liệu sản lượng và thời tiết được làm sạch tại tầng kho dữ liệu, tiến trình thực hiện tính toán biến đổi cấu trúc để chuẩn bị dữ liệu đầu vào tối ưu cho mô hình dự báo học máy.

\subsection{Giải quyết độ lệch tần suất ghi nhận (Granularity Alignment)}
Một thách thức lớn trong dữ liệu của dự án là độ lệch tần suất (Lệch pha độ phân giải thời gian) giữa hai nguồn thông tin nghiệp vụ:
\begin{itemize}
    \item Dữ liệu sản lượng điện được ghi nhận với tần suất cao (15 phút một dòng).
    \item Dữ liệu khí tượng viễn thám chỉ có sẵn theo chu kỳ giờ (60 phút một dòng).
\end{itemize}

Để giải quyết sự lệch pha này mà không làm mất mát thông tin xu hướng của bức xạ mặt trời, nhóm thiết lập thuật toán gom cụm (Aggregation) đưa dữ liệu về cùng chu kỳ giờ (1-hour grain). Chỉ số sản lượng điện mặt trời \texttt{energy\_generated\_kwh} trong 4 chu kỳ 15 phút của cùng một giờ được tính tổng cộng dồn (Sum Aggregation):
\begin{equation}
E_{\text{hour}} = \sum_{i=1}^{4} E_{15\text{min}, i}
\end{equation}
Trong khi đó, các biến thời tiết đại diện cho giờ đó được giữ nguyên giá trị trung bình giờ. Thuật toán ghép nối sử dụng hàm kết hợp chính xác mốc thời gian chẵn giờ, giúp tạo ra một bảng dữ liệu phẳng góc nhìn rộng (Wide-table) thống nhất, đảm bảo tính liên kết chặt chẽ và không có dữ liệu rò rỉ thời gian (Data Leakage).

\subsection{Trích xuất các đặc trưng thời gian và lịch trình học tập}
Điện năng mặt trời mang tính chu kỳ tự nhiên rất mạnh phụ thuộc vào quỹ đạo chuyển động của Trái Đất. Dự án tiến hành trích xuất các đặc trưng thời gian đa tầng từ trường mốc thời gian \texttt{timestamp}:
\begin{itemize}
    \item \textbf{Đặc trưng giờ trong ngày (Hour):} Giá trị từ 0 đến 23, xác định góc chiếu mặt trời.
    \item \textbf{Đặc trưng tháng trong năm (Month):} Xác định tính mùa vụ thời tiết (Mùa hè bức xạ cao, mùa đông bức xạ thấp ở Nam bán cầu).
    \item \textbf{Đặc trưng ngày trong tuần (Day of Week):} Hỗ trợ phân tích hành vi tiêu thụ điện năng tại các khu học xá (Campus).
\end{itemize}

Đặc biệt, nhóm tích hợp thành công dữ liệu lịch học tập (\texttt{calender.csv}) của cơ sở đào tạo thông qua ba đặc trưng nhị phân:
\begin{itemize}
    \item \texttt{is\_holiday}: Đánh dấu các ngày nghỉ lễ quốc gia hoặc ngày nghỉ nội bộ.
    \item \texttt{is\_semester}: Đánh dấu thời điểm đang trong học kỳ chính thức (nhu cầu sử dụng thiết bị điện cao đột biến).
    \item \texttt{is\_exam}: Đánh dấu giai đoạn thi cử tập trung.
\end{itemize}
Sự kết hợp các thuộc tính lịch trình này giúp mô hình học máy phân biệt rõ rệt giữa biến động sản lượng do thời tiết tự nhiên và biến động tiêu thụ do hành vi con người gây ra, cải thiện độ chính xác dự báo sản lượng thực dùng tại các điểm nút phụ tải trạm.

\begin{thebibliography}{99}
\addcontentsline{toc}{chapter}{Tài liệu Tham khảo}

\bibitem{wimalaratne2022unisolar}
P.~Wimalaratne, D.~Jayasinghe, và C.~Arachchige,
\newblock ``UNISOLAR: A Comprehensive High-Resolution Open Dataset of 42 Rooftop Photovoltaic Sites,''
\newblock \textit{Data in Brief}, vol.~42, p.~108152, 2022.

\bibitem{openmeteo2023}
Open-Meteo,
\newblock ``Historical Weather API \& ERA5/ERA5-Land Reanalysis Climate Dataset,''
\newblock \textit{Open-Meteo Documentation}, 2023. [Online]. Available: \url{https://open-meteo.com/}.

\bibitem{iea2023pvps}
IEA-PVPS Task 13,
\newblock ``Performance, Operation and Reliability of Photovoltaic Systems: Review of Failures and Degradation Mechanisms,''
\newblock \textit{International Energy Agency Technical Report}, IEA-PVPS T13-15:2023, 2023.

\bibitem{li2025gmmif}
Y.~Li, H.~Zhang, X.~Wang, và J.~Liu,
\newblock ``A Two-Stage Hybrid Outlier Cleaning Framework for Distributed Rooftop Photovoltaic Generation Data Using GMM and Isolation Forest,''
\newblock \textit{IEEE Access}, vol.~13, pp.~24510--24523, 2025.

\bibitem{liu2008isolation}
F.~T.~Liu, K.~M.~Ting, và Z.-H.~Zhou,
\newblock ``Isolation Forest,''
\newblock trong \textit{Proceedings of the IEEE International Conference on Data Mining (ICDM)}, 2008, pp.~413--422.

\bibitem{bishop2006pattern}
C.~M.~Bishop,
\newblock \textit{Pattern Recognition and Machine Learning},
\newblock Springer, New York, 2006.

\bibitem{tukey1977eda}
J.~W.~Tukey,
\newblock \textit{Exploratory Data Analysis},
\newblock Addison-Wesley, Reading, MA, 1977.

\bibitem{chandola2009anomaly}
V.~Chandola, A.~Banerjee, và V.~Kumar,
\newblock ``Anomaly Detection: A Survey,''
\newblock \textit{ACM Computing Surveys}, vol.~41, no.~3, pp.~1--58, 2009.

\bibitem{kim2019missing}
G.~Kim, S.~Park, và S.~Lee,
\newblock ``Physical Constraint Imputation and Spatial Interpolation for Missing Solar Irradiance Data,''
\newblock \textit{Solar Energy}, vol.~188, pp.~754--765, 2019.

\bibitem{roy2025hourly}
A.~Roy, M.~Rahman, và K.~Chowdhury,
\newblock ``Hourly-Binned Regressors Outperform Complex Deep Learning for Local Solar Generation Forecasting,''
\newblock \textit{Applied Energy}, vol.~355, p.~122301, 2025.

\bibitem{haurwitz1945}
B.~Haurwitz,
\newblock ``Insolation in Relation to Cloudiness and Cloud Density,''
\newblock \textit{Journal of Meteorology}, vol.~2, no.~3, pp.~154--166, 1945.

\bibitem{bailey2024}
J.~Bailey, R.~Miller, và S.~Thompson,
\newblock ``Clear-Sky PV Generation Modeling Under Complex Atmospheric Conditions,''
\newblock \textit{Renewable Energy}, vol.~219, pp.~119--131, 2024.

\bibitem{engerer2014kpv}
N.~A.~Engerer và F.~P.~Mills,
\newblock ``Validating the PV Clear-Sky Index and Modeling Rooftop Solar Production,''
\newblock \textit{Solar Energy}, vol.~105, pp.~619--632, 2014.

\bibitem{dev2018tes}
S.~Dev, F.~M.~Savoy, Y.~H.~Lee, và S.~Winkler,
\newblock ``High-Resolution Weather and Irradiance Monitoring for PV Plant Management,''
\newblock \textit{IEEE Transactions on Geoscience and Remote Sensing}, vol.~56, no.~4, pp.~2120--2132, 2018.

\bibitem{pin2025wape}
E.~Pin, M.~Rossi, và T.~Nguyen,
\newblock ``Weighted Absolute Percentage Error (WAPE) as a Robust Standard for Intermittent Renewable Energy Forecasting Metrics,''
\newblock \textit{Energy Reports}, vol.~11, pp.~402--415, 2025.

\bibitem{lundberg2017shap}
S.~M.~Lundberg và S.-I.~Lee,
\newblock ``A Unified Approach to Interpreting Model Predictions,''
\newblock trong \textit{Advances in Neural Information Processing Systems (NeurIPS)}, vol.~30, 2017, pp.~4765--4774.

\bibitem{zhang2013metrics}
J.~Zhang, B.-M.~Hodge, và A.~Florita,
\newblock ``Comprehensive Evaluation of Solar Irradiance and Power Forecasting Metrics,''
\newblock \textit{Journal of Renewable and Sustainable Energy}, vol.~5, no.~4, p.~043126, 2013.

\bibitem{lonij2012tucson}
V.~P.~Lonij, A.~E.~Brooks, A.~D.~Cronin, và F.~J.~Rocha,
\newblock ``Intra-Hour Solar Forecasting with High-Density PV Networks and Cloud Tracking,''
\newblock \textit{Solar Energy}, vol.~86, no.~5, pp.~1450--1461, 2012.

\bibitem{taylor2018prophet}
S.~J.~Taylor và B.~Letham,
\newblock ``Forecasting at Scale,''
\newblock \textit{The American Statistician}, vol.~72, no.~1, pp.~37--45, 2018.

\bibitem{box2015time}
G.~E.~P.~Box, G.~M.~Jenkins, G.~C.~Reinsel, và G.~M.~Ljung,
\newblock \textit{Time Series Analysis: Forecasting and Control}, 5th ed.,
\newblock John Wiley \& Sons, Hoboken, NJ, 2015.

\bibitem{ke2017lightgbm}
G.~Ke, Q.~Meng, T.~Finley, T.~Wang, W.~Chen, W.~Ma, Q.~Ye, và T.-Y.~Liu,
\newblock ``LightGBM: A Highly Efficient Gradient Boosting Decision Tree,''
\newblock trong \textit{Advances in Neural Information Processing Systems (NeurIPS)}, vol.~30, 2017, pp.~3146--3154.

\bibitem{chen2016xgboost}
T.~Chen và C.~Guestrin,
\newblock ``XGBoost: A Scalable Tree Boosting System,''
\newblock trong \textit{Proceedings of the ACM SIGKDD International Conference on Knowledge Discovery and Data Mining}, 2016, pp.~785--794.

\bibitem{kimball2013dwh}
R.~Kimball và M.~Ross,
\newblock \textit{The Data Warehouse Toolkit: The Definitive Guide to Dimensional Modeling}, 3rd ed.,
\newblock John Wiley \& Sons, Indianapolis, IN, 2013.

\bibitem{wmo2021guide}
World Meteorological Organization (WMO),
\newblock \textit{Guide to Instruments and Methods of Observation: Volume I -- Measurement of Meteorological Variables},
\newblock WMO-No. 8, Geneva, Switzerland, 2021.

\bibitem{supabase2024}
Supabase Inc.,
\newblock ``PostgreSQL Database \& Row-Level Security for Modern Analytics Architecture,''
\newblock \textit{Supabase Technical Architecture Documentation}, 2024. [Online]. Available: \url{https://supabase.com/docs}.

\bibitem{pedregosa2011sklearn}
F.~Pedregosa et~al.,
\newblock ``Scikit-learn: Machine Learning in Python,''
\newblock \textit{Journal of Machine Learning Research}, vol.~12, pp.~2825--2830, 2011.

\bibitem{tableau2023guidelines}
Tableau Software (Salesforce),
\newblock \textit{Visual Analysis Best Practices: A Guide to Creating Meaningful Dashboards},
\newblock Seattle, WA, 2023.

\end{thebibliography}


\end{document}
